\documentclass{article}
\usepackage[utf8]{inputenc}
\usepackage{amsmath,amssymb}
\usepackage[most]{tcolorbox}
\usepackage{geometry}
\geometry{margin=2.5cm}

\newcommand{\step}[1]{\par\noindent\hangindent=3.5em\hangafter=1%
  \makebox[3.5em][l]{\textbf{#1.}}}
\newcommand{\steptag}[1]{\hfill{\small\textsf{[#1]}}}

\newtcolorbox{proofbox}[1]{
  colback=gray!5, colframe=gray!60,
  fonttitle=\bfseries, title={#1},
  breakable, enhanced,
  left=4pt, right=4pt, top=4pt, bottom=4pt,
  before skip=10pt, after skip=10pt
}

\begin{document}

\begin{proofbox}{After first fixing one global witness package W\_RF suppli...}
\step{1} After first fixing one global witness package W\_RF supplied by lem-routef-raw-factor-setting-formation, for every input (H,Phi,eta) to which that formation result applies, fix one def-routef-raw-factor-setting datum S over that same W\_RF supplied by the same result, writing the fields of (W\_RF,S) as the unqualified symbols below: for every Delta' supplied for (W\_RF,S) by lem-routef-delta-prime-closeness, every Delta supplied from that same Delta' by lem-routef-delta-normalization-closeness, and every componentwise package (m,(L\_j,E\_j,W\_j,Sigma\_j,U\_js,p\_js,C\_j,xi\_j,Lambda\_j,Upsilon'\_j)\_\{j=1\}\textasciicircum{}m,F,V,Upsilon') supplied for (W\_RF,S,Delta',Delta) by lem-routef-upsilon-prime-component-construction, with C\_L:=C\_2+C\_3+2*C\_R from (1.3) and rho\_Upsilon' := min\{rho\_T, rho\_id, rho\_Delta, rho\_2, rho\_3, (2*C\_R)\textasciicircum{}(-1)\}, for 0 <= eta <= rho\_Upsilon', every integer q >= 1, and every Y=(Y\_1,...,Y\_m) in M\_q(B)=direct-sum\_\{j=1\}\textasciicircum{}m M\_q(B(L\_j)), ||(Upsilon'\_j)\_q(Delta\_q(Y))-Y\_j|| <= C\_L*eta*||Y|| for every j, and consequently ||(Upsilon' Delta-I\_B)\_q(Y)|| <= C\_L*eta*||Y|| and ||Upsilon' Delta-I\_B||\_cb <= C\_L*eta. \steptag{validated} \\[6pt]
\step{1.1} Fix the data and hypotheses of node 1, fix q>=1, Y in M\_q(B), and j. Put P\_jq(Y):=sum\_\{s,t in Sigma\_j\} p\_js p\_jt iota\_js,q\textasciicircum{}* W\_j,q Delta\_q(hat-U\_js Y hat-U\_jt\textasciicircum{}*) W\_j,q\textasciicircum{}* iota\_jt,q. Then ||(Upsilon'\_j)\_q(Delta\_q(Y))-P\_jq(Y)|| <= (C\_2+C\_3)*eta*||Y||. \steptag{validated} \\[6pt]
\step{1.1.1} Because rho\_Upsilon'=min\{rho\_T,rho\_id,rho\_Delta,rho\_2,rho\_3,(2*C\_R)\textasciicircum{}(-1)\}, 0<=eta<=rho\_Upsilon' implies eta<=rho\_Delta,rho\_2,rho\_3. By lem-routef-delta-normalization-closeness Delta is UCP, while Phi is UCP in the def-routef-raw-factor-setting datum supplied by lem-routef-raw-factor-setting-formation. Hence every Delta\_q and Phi\_q is unital and contractive. Also lem-routef-upsilon-prime-component-construction gives sum\_s p\_js=1 with p\_js>=0, sum\_j W\_j\textasciicircum{}*W\_j=I, unit xi\_j and unitary U\_js; consequently ||W\_j,q||<=1, each iota\_js,q is an isometry, and ||hat-U\_js||=1. \steptag{validated} \\[4pt]
\step{1.1.2} Let A\_st:=Phi\_q(Delta\_q(hat-U\_js) Phi\_q(Delta\_q(Y)) Delta\_q(hat-U\_jt\textasciicircum{}*)). Expanding Lambda\_j and Upsilon'\_j from lem-routef-upsilon-prime-component-construction at level q, and using Phi\_q(T)=V\_q\textasciicircum{}*(T tensor I\_F)V\_q, gives the exact identity (Upsilon'\_j)\_q(Delta\_q(Y))=sum\_\{s,t\}p\_js p\_jt iota\_js,q\textasciicircum{}* W\_j,q A\_st W\_j,q\textasciicircum{}* iota\_jt,q. \steptag{validated} \\[4pt]
\step{1.1.3} Applying lem-routef-degree-two-estimate at n=q with X=I and the present Y, and using Delta\_q(I)=I, gives ||Phi\_q(Delta\_q(Y))-Delta\_q(Y)||<=C\_2*eta*||Y||. Contractivity of Phi\_q and Delta\_q and ||hat-U\_js||=||hat-U\_jt\textasciicircum{}*||=1 therefore imply ||A\_st-Phi\_q(Delta\_q(hat-U\_js)Delta\_q(Y)Delta\_q(hat-U\_jt\textasciicircum{}*))||<=C\_2*eta*||Y||. \steptag{validated} \\[6pt]
\step{1.1.3.1} For fixed s,t in Sigma\_j, define A\_st exactly as in node 1.1.2 by A\_st:=Phi\_q(Delta\_q(hat-U\_js) Phi\_q(Delta\_q(Y)) Delta\_q(hat-U\_jt\textasciicircum{}*)). By node 1.1.1, Delta\_q(I)=I and Phi\_q,Delta\_q are contractive. Applying lem-routef-degree-two-estimate at n=q with X=I and the present Y gives ||Phi\_q(Delta\_q(Y))-Delta\_q(Y)||<=C\_2*eta*||Y||. With D:=Phi\_q(Delta\_q(Y))-Delta\_q(Y), linearity gives A\_st-Phi\_q(Delta\_q(hat-U\_js)Delta\_q(Y)Delta\_q(hat-U\_jt\textasciicircum{}*))=Phi\_q(Delta\_q(hat-U\_js) D Delta\_q(hat-U\_jt\textasciicircum{}*)). Hence contractivity and ||hat-U\_js||=||hat-U\_jt\textasciicircum{}*||=1 imply the asserted bound <=C\_2*eta*||Y||. \steptag{validated} \\[4pt]
\step{1.1.4} By lem-routef-degree-three-estimate at n=q with X=hat-U\_js, Y=Y, Z=hat-U\_jt\textasciicircum{}*, the second term in node 1.1.3 differs from Delta\_q(hat-U\_js Y hat-U\_jt\textasciicircum{}*) by at most C\_3*eta*||Y||. Thus each A\_st has error at most (C\_2+C\_3)*eta*||Y||. Sandwiching by the contractions iota\_js,q\textasciicircum{}* W\_j,q and W\_j,q\textasciicircum{}* iota\_jt,q, summing with nonnegative weights, and using sum\_\{s,t\}p\_js p\_jt=1 together with node 1.1.2 proves node 1.1. \steptag{validated} \\[4pt]
\step{1.2} Under the same fixed data, define P\_jq(Y):=sum\_\{s,t in Sigma\_j\} p\_js p\_jt iota\_js,q\textasciicircum{}* W\_j,q Delta\_q(hat-U\_js Y hat-U\_jt\textasciicircum{}*) W\_j,q\textasciicircum{}* iota\_jt,q. Then P\_jq(Y)=alpha\_j Y\_j where alpha\_j:=<xi\_j,C\_j\textasciicircum{}*C\_j xi\_j>, and ||P\_jq(Y)-Y\_j|| <= 2*C\_R*eta*||Y||. \steptag{validated} \\[6pt]
\step{1.2.1} Let P\_j:=W\_jW\_j\textasciicircum{}*, let iota\_j0(z):=z tensor xi\_j, let u\_js,q:=I\_q tensor U\_js on C\textasciicircum{}q tensor L\_j, and let Ucal\_js,q:=I\_q tensor U\_js tensor I\_Ej on C\textasciicircum{}q tensor L\_j tensor E\_j. From Delta(X)=sum\_k W\_k\textasciicircum{}*(X\_k tensor I\_Ek)W\_k in lem-routef-upsilon-prime-component-construction and the fact that hat-U\_js Y hat-U\_jt\textasciicircum{}* has only block j nonzero, Delta\_q(hat-U\_js Y hat-U\_jt\textasciicircum{}*)=W\_j,q\textasciicircum{}*[(u\_js,q Y\_j u\_jt,q\textasciicircum{}*) tensor I\_Ej]W\_j,q. Substitution in P\_jq(Y), with iota\_js,q=Ucal\_js,q iota\_j0,q, gives P\_jq(Y)=iota\_j0,q\textasciicircum{}* [sum\_s p\_js Ucal\_js,q\textasciicircum{}* P\_j,q Ucal\_js,q] (Y\_j tensor I\_Ej) [sum\_t p\_jt Ucal\_jt,q\textasciicircum{}* P\_j,q Ucal\_jt,q] iota\_j0,q. \steptag{validated} \\[4pt]
\step{1.2.2} By the identity R\_j=sum\_s p\_js (U\_js\textasciicircum{}* tensor I\_Ej)P\_j(U\_js tensor I\_Ej)=I\_Lj tensor C\_j in lem-routef-upsilon-prime-component-construction, both bracketed sums in node 1.2.1 equal I\_q tensor R\_j=I\_q tensor I\_Lj tensor C\_j. Therefore P\_jq(Y)=iota\_j0,q\textasciicircum{}*(I\_q tensor I\_Lj tensor C\_j)(Y\_j tensor I\_Ej)(I\_q tensor I\_Lj tensor C\_j)iota\_j0,q=<xi\_j,C\_j\textasciicircum{}*C\_j xi\_j>Y\_j=alpha\_jY\_j. \steptag{validated} \\[4pt]
\step{1.2.3} Since C\_j is a positive contraction and xi\_j is a unit vector, 0<=alpha\_j=<xi\_j,C\_j\textasciicircum{}*C\_j xi\_j><=1, so |1-alpha\_j|=1-alpha\_j. The component-construction bound 1-<xi\_j,C\_j\textasciicircum{}*C\_j xi\_j><=2*C\_R*eta and ||Y\_j||<=||Y|| now give ||P\_jq(Y)-Y\_j||=|1-alpha\_j|*||Y\_j||<=2*C\_R*eta*||Y||, proving node 1.2. \steptag{validated} \\[4pt]
\step{1.2.4} The opening clause of node 1.2 now binds P\_jq(Y) explicitly by the same finite double-sum formula used in node 1.1. Therefore every occurrence of P\_jq(Y) in nodes 1.2.1--1.2.3 is in scope from its ancestor 1.2, and node 1.3 combines the estimates of nodes 1.1 and 1.2 for one and the same explicitly defined operator P\_jq(Y). \steptag{validated} \\[4pt]
\step{1.3} By the two preceding estimates and C\_L=C\_2+C\_3+2*C\_R, every j satisfies ||(Upsilon'\_j)\_q(Delta\_q(Y))-Y\_j|| <= C\_L*eta*||Y||. \steptag{validated} \\[4pt]
\step{1.4} Since Upsilon'(T)=(Upsilon'\_j(T))\_j and M\_q(B) has the finite direct-sum maximum norm, the component estimates give ||(Upsilon' Delta-I\_B)\_q(Y)|| <= C\_L*eta*||Y|| for every q,Y; taking the supremum over q and nonzero Y gives ||Upsilon' Delta-I\_B||\_cb <= C\_L*eta. \steptag{validated} \\[4pt]
\end{proofbox}

\end{document}
